\renewcommand{\thestatement}{4.\number\numexpr\value{statement}-78\relax}

\subsection{Laplace 问题}
\label{sec:chap3-laplace-problem}

本节先研究 Dirichlet 边界条件, 再研究 Neumann 边界条件. 存在唯一性由 Lax--Milgram 定理和迹提升得到; 正则性则用前面建立的 Nirenberg 平移方法自足地证明.

\subsubsection{Dirichlet 边界条件}
\label{subsec:chap3-laplace-dirichlet}

\begin{Theorem}
\label{thm:chap3-laplace-dirichlet-wellposedness}
设 $\Omega\subset\mathbb R^d$ 为有界连通 Lipschitz 区域, $f\in H^{-1}(\Omega)$, $u_b\in H^{1/2}(\partial\Omega)$. 则问题
\begin{equation}
 \begin{cases}
 -\Delta u=f,&\text{于 }\Omega,\\
 u=u_b,&\text{于 }\partial\Omega
 \end{cases}
 \label{eq:chap3-laplace-dirichlet-problem}
\end{equation}
存在唯一解 $u\in H^1(\Omega)$, 且
\[
 \|u\|_{H^1}\leq C(\|f\|_{H^{-1}}+\|u_b\|_{H^{1/2}}).
\]
\end{Theorem}

\begin{Proof}
取 Theorem~\ref{thm:chap3-trace-lifting-operator} 的提升 $R_0u_b$. 则 $\widetilde u=u-R_0u_b\in H_0^1(\Omega)$ 满足齐次边界问题, 右端为 $\widetilde f=f+\Delta(R_0u_b)\in H^{-1}$. 在 $H_0^1(\Omega)$ 上定义
\[
 a(u,v)=\int_\Omega\nabla u\cdot\nabla v\,\mathrm dx.
\]
Cauchy--Schwarz 不等式给出连续性, Proposition~\ref{prop:chap3-generalized-poincare} 给出强制性. Lax--Milgram 定理产生唯一 $\widetilde u$, 并给出所需估计. 对 $\mathcal D(\Omega)$ 测试即恢复分布方程, 而 Theorem~\ref{thm:chap3-zero-trace-characterization}给出边界条件.
\end{Proof}

\begin{Theorem}
\label{thm:chap3-laplace-dirichlet-regularity}
沿用 Theorem~\ref{thm:chap3-laplace-dirichlet-wellposedness} 的记号. 若 $k\geq0$, $\Omega$ 为 $C^{k+1,1}$ 区域, $f\in H^k(\Omega)$, $u_b\in H^{k+3/2}(\partial\Omega)$, 则式\eqref{eq:chap3-laplace-dirichlet-problem}的解属于 $H^{k+2}(\Omega)$, 且
\[
 \|u\|_{H^{k+2}}\leq C(\|f\|_{H^k}+\|u_b\|_{H^{k+3/2}}).
\]
\end{Theorem}

\begin{Proof}
只详细处理 $k=0$; 高阶情形逐次微分重复同一论证. Theorem~\ref{thm:chap3-trace-lifting-operator} 的高阶性质把非齐次边界数据化为齐次情形.

先证内部正则性. 对 $\varphi\in\mathcal D(\Omega)$, 令 $v=u\varphi$ 并在全空间计算
\[
 -\Delta v=f\varphi-2\nabla u\cdot\nabla\varphi-u\Delta\varphi=:F\in L^2(\mathbb R^d).
\]
于是 $-\Delta(\delta_hv)=\delta_hF$. 以 $\delta_hv$ 测试并应用 Lemma~\ref{lem:chap3-translation-sobolev-bound}与 Young 不等式, 得 $\|\nabla\delta_hv\|_{L^2}\leq C|h|$. Proposition~\ref{prop:chap3-sobolev-difference-quotient-characterization}推出 $v\in H^2(\mathbb R^d)$, 故 $u\in H^2_{\mathrm{loc}}(\Omega)$.

再证切向正则性. 齐次问题的变分式为
\begin{equation}
 \int_\Omega\nabla u\cdot\nabla v\,\mathrm dx
 =\int_\Omega fv\,\mathrm dx,
 \qquad v\in H_0^1(\Omega).
 \label{eq:chap3-dirichlet-variational-form}
\end{equation}
对满足式\eqref{eq:chap3-admissible-tangential-field}的 $\theta$, 取弯曲差商测试函数 $\delta_{-h}^\theta\delta_h^\theta u$. Proposition~\ref{prop:chap3-curved-translation-properties}和强制性给出
\[
 \|\delta_h^\theta u\|_{H_0^1}\leq C|h|(\|f\|_{L^2}+\|\nabla u\|_{L^2}).
\]
Theorem~\ref{thm:chap3-tangential-space-characterization}推出 $u\in H_{mathrm{tang}}^2(\Omega)$.

最后在边界条带上使用 Proposition~\ref{prop:chap3-laplacian-normal-tangential}:
\begin{equation}
 \operatorname{Tr}(\nabla_T\nabla u)+\nabla_N\nabla_Nu
 +\nabla u\cdot\widetilde G=f.
 \label{eq:chap3-laplace-normal-tangential-equation}
\end{equation}
几何系数满足
\begin{equation}
 \frac{\widetilde G}{|\ln\rho|}\in L^\infty(O_{s_0}^\rho)^d,
 \label{eq:chap3-laplace-geometric-coefficient-bound}
\end{equation}
故先有
\begin{equation}
 \nabla u\cdot\widetilde G\in L^q(O_{s_0}^\rho),
 \qquad\forall q<2.
 \label{eq:chap3-laplace-geometric-product-integrability}
\end{equation}
式\eqref{eq:chap3-laplace-normal-tangential-equation}于是给出 $u\in W^{2,q}$ 对每个 $q<2$. Sobolev 嵌入再提高 $\nabla u$ 的可积指数至某个 $r>2$, 从而式\eqref{eq:chap3-laplace-geometric-product-integrability}可取 $q=2$. 回代即得 $u\in H^2(\Omega)$ 及估计.
\end{Proof}

\subsubsection{Neumann 边界条件}
\label{subsec:chap3-laplace-neumann}

\begin{Theorem}
\label{thm:chap3-laplace-neumann}
设 $\Omega$ 为有界连通 Lipschitz 区域, $f\in L^2(\Omega)$, $f_b\in H^{-1/2}(\partial\Omega)$, 且满足相容条件
\begin{equation}
 \int_\Omega f(x)\,\mathrm dx
 +\langle f_b,1\rangle_{H^{-1/2},H^{1/2}}=0.
 \label{eq:chap3-neumann-compatibility}
\end{equation}
则存在唯一的 $u\in H^1(\Omega)\cap L_0^2(\Omega)$ 满足
\[
 -\Delta u=f\quad\text{于 }\Omega,
 \qquad\partial_\nu u=f_b\quad\text{于 }\partial\Omega,
\]
并且
\begin{equation}
 \|u\|_{H^1}\leq C(\|f\|_{L^2}+\|f_b\|_{H^{-1/2}}).
 \label{eq:chap3-neumann-energy-bound}
\end{equation}
若进一步 $\Omega$ 为 $C^{k+1,1}$, $f\in H^k$, $f_b\in H^{k+1/2}$, 则 $u\in H^{k+2}$ 且
\[
 \|u\|_{H^{k+2}}\leq C(\|f\|_{H^k}+\|f_b\|_{H^{k+1/2}}).
\]
\end{Theorem}

\textbf{Remark 4.1:} 若 $u\in H^1$ 且 $-\Delta u=f$ 于分布意义成立, 则 $\nabla u$ 属于后文研究的 $H_{\mathrm{div}}$ 空间, 因而法向导数在 $H^{-1/2}(\partial\Omega)$ 中有弱意义. 广义 Stokes 公式表明式\eqref{eq:chap3-neumann-compatibility}是可解性的必要条件; 零均值则保证唯一性.

\begin{Proof}
令
\[
 H_m^1(\Omega)=\left\{v\in H^1(\Omega):\int_\Omega v\,\mathrm dx=0\right\},
\]
并在其上定义
\[
 a(u,v)=\int_\Omega\nabla u\cdot\nabla v\,\mathrm dx,
 \quad L(v)=\int_\Omega fv\,\mathrm dx+langle f_b,\gamma_0v\rangle.
\]
Proposition~\ref{prop:chap3-poincare-wirtinger}使 $a$ 强制, 故 Lax--Milgram 定理给出唯一零均值解和式\eqref{eq:chap3-neumann-energy-bound}. 式\eqref{eq:chap3-neumann-compatibility}说明 $L(1)=0$, 因而变分式对所有 $H^1$ 测试函数成立, 这同时给出方程与弱 Neumann 条件.

正则性方面, 先用 Theorem~\ref{thm:chap3-normal-trace-lifting}提升 $f_b$, 化为齐次 Neumann 条件. 内部 $H^2$ 正则性与 Dirichlet 情形完全相同. 切向差商不再具有零迹, 但 Proposition~\ref{prop:chap3-curved-translation-properties}和 $\int_\Omega u=0$ 给出
\[
 \left|\int_\Omega\delta_h^\theta u\,\mathrm dx\right|
 \leq C|h|\|u\|_{L^2}.
\]
用 Poincar\'e--Wirtinger 不等式替代齐次 Poincar\'e 不等式, 可重复切向正则性论证. 最后通过式\eqref{eq:chap3-laplace-normal-tangential-equation}恢复完整正则性, 与边界条件类型无关.
\end{Proof}
